\section{Discussion}
\label{sec:discussion}

The central result is not simply that two national totals agree. In the latest two fiscal
years, OCAC commercial fuel and OECD combined territory aviation differ by less than 1\%, but
resident and foreign components differ by about $+29$\% and $-23$\%. The signs persist in the
monthly record, and the five-year combined ratio ranges from 0.892 to 0.998. Aggregate closure
is therefore a useful check on the gross boundary, but it cannot establish that customer
categories have been assigned correctly. Indeed, the closer the cancellation, the easier it
is for a component error to disappear from the national total.

The conserved ledger changes how that evidence is interpreted. It prevents a customer label
from being promoted directly to an IPCC category, preserves the international counterpart of
every domestic assignment and leaves unsupported fuel as a statistical discrepancy. It also
separates four quantities that serve different users: domestic aviation under the IPCC flight
boundary, full-flight international bunkers, resident-carrier emissions and emissions located
physically over Pakistan. A single fuel-sales label cannot represent all four.

Pakistan's official submission demonstrates the policy relevance. Its 80/20 distribution
produces domestic and international estimates on opposite sides of the flight-geography
evidence by factors of 2.81 and 0.095. The official values should not be dismissed merely
because they differ from OECD, which is itself modelled. However, the aviation-gasoline label,
the absence of estimated jet kerosene and the inability to reconstruct the allocation from
public documentation mean that the split warrants institutional review. The appropriate next
step is to reconcile the CRT activity data with OCAC jet-fuel rows and actual domestic and
international flight stages, rather than replace one opaque percentage with another.

The hierarchy of domestic estimates localises the remaining scientific uncertainty. Across the
complete inventory, replacing the 550-nmi representative stage with the constrained route
matrix lowers fuel by 9.4\%. On the identical 17,491-stage subset used for the resolution
experiment, the corresponding reduction is 6.6\%. Adding aircraft-specific OpenAP profiles on
that common support lowers fuel by a further 3.5\% after a physical range check corrected the
earlier software interface error. The route--aircraft Monte Carlo mean of 0.244 Mt \COtwo{} is
close to OECD's 0.236 Mt, but this is model-to-model agreement, not observed block-fuel
validation. Its
probabilistic and structural ranges identify aircraft assignment, detour and completion rules
as actionable uncertainties. PCAA passenger, cargo and mail data already constrain payload;
registrations, route-complete international activity and airline block fuel would test the
remaining assumptions directly.

The external comparisons have deliberately different roles. Climate TRACE reproduces the
airport ranking and broadly agrees on fiscal-year magnitude, but may share flight-surveillance
and performance inputs with OECD or a future ADS-B implementation. The negative ADSB.lol pilot
does not undermine the accounting result: positive controls show that the adapter works, while
the Pakistan sample does not meet the coverage gate. Observed flight records would replace one
input tier and test the present estimates; they are not required to make source-row conservation
or component cancellation valid.

The approach is transferable wherever petroleum statistics follow purchasers while greenhouse-
gas inventories follow flight geography. The essential requirements are modest: preserve each
fuel source row, distinguish customer from operator and endpoints, disclose the role and shared
inputs of every comparison, and report structural mappings separately from parametric
uncertainty. More detailed trajectory or aircraft models then improve a defined layer of the
account rather than silently changing its boundary.
